\section{Conclusion}
\label{sect:conclusion}

This paper provides the first system-level characterization of AI agents from an AI infrastructure perspective. While these LLM-based agents demonstrate powerful reasoning capabilities, they also introduce substantial energy overheads that are orders of magnitude higher than conventional single-turn LLM inference. Our analysis shows that common agent design patterns
incur heavy latency penalties and infrastructure costs, especially when deployed at scale. Moreover, test-time scaling yields sharply diminishing returns in accuracy, challenging the cost-effectiveness of current agent implementations.

These findings underscore an urgent need to rethink agent architecture and workflow design. Rather than relying on brute-force test-time scaling, future agents should adopt compute-aware reasoning strategies that optimize accuracy per unit cost. This includes smarter scheduling, caching, prompt engineering, and hybrid scaling approaches that adapt to deployment constraints. By exposing the hidden costs of agentic reasoning and offering actionable insights into their infrastructure impact, we hope this work informs future system and algorithm co-design for scalable and sustainable AI agents.

